\documentclass[10pt, aspectratio=169]{beamer}
\usetheme{Madrid}
\usecolortheme{dolphin}
\usepackage{ctex}
\usepackage{xeCJK}
\usepackage{amsmath,amssymb}
\usepackage{listings}
\usepackage{tikz}
\usepackage{xcolor}

\title{408考研零基础 --- 导学篇}
\subtitle{0-4 C语言速成（上）}
\author{}
\date{}
\setbeamertemplate{page number in head/foot}{}

\begin{document}

\begin{frame}{本节目标}
    \begin{enumerate}
        \item 掌握 C 语言的基本语法结构
        \item 理解变量、常量的概念与基本数据类型
        \item 掌握三大控制结构：顺序、选择、循环
    \end{enumerate}
\end{frame}

\begin{frame}{C 语言程序结构}
    \begin{lstlisting}[language=C, basicstyle=\ttfamily\small]
#include <stdio.h>

int main() {
    printf("Hello, World!\n");
    return 0;
}
    \end{lstlisting}
    \vspace{0.3em}
    \textbf{组成部分}
    \begin{itemize}
        \item 预处理指令：\#include
        \item 主函数：int main()
        \item 函数体：变量定义 + 语句
        \item return 语句退出程序
    \end{itemize}
\end{frame}

\begin{frame}{基本数据类型}
    \begin{center}
    \begin{tabular}{|c|c|c|}
        \hline
        \textbf{类型} & \textbf{关键字} & \textbf{字节} \\ \hline
        整型 & int & 4 \\ \hline
        字符型 & char & 1 \\ \hline
        单精度浮点 & float & 4 \\ \hline
        双精度浮点 & double & 8 \\ \hline
        无值型 & void & --- \\ \hline
    \end{tabular}
    \end{center}
    \vspace{0.5em}
    \textbf{声明示例}
    \begin{lstlisting}[language=C, basicstyle=\ttfamily\small]
int age = 20;
char grade = 'A';
double price = 99.99;
const double PI = 3.14159;
    \end{lstlisting}
\end{frame}

\begin{frame}{运算符}
    \textbf{算术运算符}
    \[
    +,\; -,\; \times,\; \div,\; \%
    \]
    \vspace{0.3em}
    \textbf{自增自减}
    \[
    i++ \text{（先使用后自增）},\; ++i \text{（先自增后使用）}
    \]
    \vspace{0.3em}
    \textbf{关系运算符}
    \[
    >,\; <,\; >=\; <=\; ==\; !=
    \]
    \vspace{0.3em}
    \textbf{逻辑运算符}
    \[
    \&\& \text{（与）},\; || \text{（或）},\; ! \text{（非）}
    \]
\end{frame}

\begin{frame}{顺序结构}
    \textbf{printf --- 格式化输出}
    \begin{lstlisting}[language=C, basicstyle=\ttfamily\small]
int a = 10;
printf("a = %d\n", a);
    \end{lstlisting}
    \vspace{0.3em}
    \textbf{scanf --- 格式化输入}
    \begin{lstlisting}[language=C, basicstyle=\ttfamily\small]
int b;
scanf("%d", &b);
    \end{lstlisting}
    \vspace{0.3em}
    \begin{tabular}{|c|c|}
        \hline \textbf{格式符} & \textbf{含义} \\ \hline
        \%d & 整型 \\ \hline
        \%c & 字符型 \\ \hline
        \%f & 浮点型 \\ \hline
    \end{tabular}
\end{frame}

\begin{frame}[fragile]{选择结构 --- if}
    \begin{lstlisting}[language=C, basicstyle=\ttfamily\small]
if (score >= 60) {
    printf("Pass\n");
} else {
    printf("Fail\n");
}
    \end{lstlisting}
    \vspace{0.3em}
    \textbf{多分支}
    \begin{lstlisting}[language=C, basicstyle=\ttfamily\small]
if (score >= 90)      grade = 'A';
else if (score >= 80) grade = 'B';
else if (score >= 70) grade = 'C';
else if (score >= 60) grade = 'D';
else                  grade = 'F';
    \end{lstlisting}
\end{frame}

\begin{frame}[fragile]{选择结构 --- switch}
    \begin{lstlisting}[language=C, basicstyle=\ttfamily\small]
switch (op) {
    case '+': result = a + b; break;
    case '-': result = a - b; break;
    case '*': result = a * b; break;
    case '/': result = a / b; break;
    default: printf("Invalid\n");
}
    \end{lstlisting}
    \vspace{0.3em}
    注意：每个 case 后需要 break，否则会发生穿透
\end{frame}

\begin{frame}[fragile]{循环结构}
    \textbf{for 循环}
    \begin{lstlisting}[language=C, basicstyle=\ttfamily\small]
for (int i = 0; i < 10; i++) {
    printf("%d ", i);
}
    \end{lstlisting}
    \vspace{0.3em}
    \textbf{while 循环}
    \begin{lstlisting}[language=C, basicstyle=\ttfamily\small]
int i = 0;
while (i < 10) {
    printf("%d ", i++);
}
    \end{lstlisting}
\end{frame}

\begin{frame}[fragile]{do-while 循环与循环控制}
    \textbf{do-while}
    \begin{lstlisting}[language=C, basicstyle=\ttfamily\small]
int i = 0;
do {
    printf("%d ", i++);
} while (i < 10);
    \end{lstlisting}
    \vspace{0.3em}
    \textbf{break / continue}
    \begin{itemize}
        \item break --- 跳出当前循环
        \item continue --- 跳过本次循环剩余部分
    \end{itemize}
\end{frame}

\begin{frame}{本节总结}
    \begin{enumerate}
        \item C 程序结构：预处理 $\to$ main 函数 $\to$ 函数体
        \item 数据类型：int, char, float, double, void
        \item 控制结构：顺序、选择 (if/switch)、循环 (for/while/do-while)
        \item C 语言是数据结构算法描述的主要语言
    \end{enumerate}
    \vspace{1em}
    \begin{center}
        \textcolor{blue}{\textbf{导学篇结束，下节开始数据结构模块！}}
    \end{center}
\end{frame}

\end{document}
